% id: 34-05
% book: 베이직쎈 미적분1 · 02 함수의 연속
% section: 유형 014 함수의 그래프와 연속
% figure: tikz:fig-34-05
% answer: 
% answer_source: 
% variant_level: 0 (원본 전사)
% 이 파일은 build.mjs 가 items.json 에서 만든다. 고칠 때는 items.json 을 고친다.
\smash{\raisebox{1.5mm}{\dmkichul{베이직쎈 미적분1 34쪽 05번}}}
\noindent\begin{minipage}[t]{0.58\linewidth}\vspace{0pt}\raggedright 구간 $(0{,}\mkern3mu\mathopen{}4)$에서 함수 $y=f(x)$의 그래프가 오른쪽 그림과 같을 때, 옳은 것만을 \textbf{보기}에서 있는 대로 고른 것은?\end{minipage}\hfill\begin{minipage}[t]{0.4\linewidth}\vspace{0pt}\centering \resizebox{\ifdim\width>\linewidth\linewidth\else\width\fi}{!}{% 34-05(34쪽) — 구간 $(0,\,4)$ 에서 $y=f(x)$ 의 그래프. 스캔 화질이 낮아 TikZ 로 다시 그림.
% 원문에 있는 것만: 빈 점 $(0,\,0)$·$(1,\,1)$·$(3,\,2)$·$(4,\,3)$ · 채운 점 $(3,\,3)$ · 점선 · 라벨 $1$, $2$, $3$, $4$, O, x, y, $y=f(x)$(눈금은 만들지 않는다).
\begin{tikzpicture}[x=0.62cm, y=0.62cm, line width=0.6pt]
  \draw[->, >=stealth, line width=0.7pt] (-0.55,0) -- (5.05,0) node[below=1pt, font=\small] {$x$};
  \draw[->, >=stealth, line width=0.7pt] (0,-0.9) -- (0,4.05) node[left=1pt, font=\small] {$y$};
  \draw[densely dashed, line width=0.4pt] (0,1) -- (1,1);
  \draw[densely dashed, line width=0.4pt] (0,2) -- (3,2);
  \draw[densely dashed, line width=0.4pt] (0,3) -- (3,3);
  \draw[densely dashed, line width=0.4pt] (1,0) -- (1,1);
  \draw[densely dashed, line width=0.4pt] (2,0) -- (2,3);
  \draw[densely dashed, line width=0.4pt] (3,0) -- (3,3);
  \draw[densely dashed, line width=0.4pt] (4,0) -- (4,3);
  \draw (0,0) -- (1,1);
  \draw[domain=1:2, samples=40, smooth] plot (\x, {3-2*(2-\x)*(2-\x)});
  \draw[domain=2:3, samples=40, smooth] plot (\x, {3-(\x-2)*(\x-2)});
  \draw (3,3) -- (4,3);
  \draw[fill=white] (0,0) circle (2.2pt);
  \draw[fill=white] (1,1) circle (2.2pt);
  \draw[fill=white] (3,2) circle (2.2pt);
  \fill (3,3) circle (2.2pt);
  \draw[fill=white] (4,3) circle (2.2pt);
  \node[below=1pt, font=\small] at (1,0) {$1$};
  \node[below=1pt, font=\small] at (2,0) {$2$};
  \node[below=1pt, font=\small] at (3,0) {$3$};
  \node[below=1pt, font=\small] at (4,0) {$4$};
  \node[below left=-2pt and -4pt, font=\small] at (0,0) {$\pt{O}$};
  \node[left=2pt, font=\small] at (0,1) {$1$};
  \node[left=2pt, font=\small] at (0,2) {$2$};
  \node[left=2pt, font=\small] at (0,3) {$3$};
  \node[anchor=west, font=\small] at (2.55,3.75) {$y=f(x)$};
\end{tikzpicture}
}\end{minipage}\par
\noindent \bogi{ㄱ. $\displaystyle\lim_{x\to 1-} f(x)=1$\\ ㄴ. $\displaystyle\lim_{x\to 3} f(x)$의 값이 존재한다.\\ ㄷ. 함수 $f(x)$는 $x=2$에서 연속이다.\\ ㄹ. 구간 $(0{,}\mkern3mu\mathopen{}4)$에서 함수 $f(x)$가 불연속인 모든 $x$의 값의 합은 6이다.}
\dmchoicesauto{}{ㄱ, ㄴ}{ㄱ, ㄷ}{ㄴ, ㄷ}{ㄱ, ㄷ, ㄹ}{ㄴ, ㄷ, ㄹ}
